\chapter{Tests, résultats et validation}
\label{chap:tests-conclusion}

\section{Introduction}
\label{sec:tc-intro}

Ce chapitre confronte la solution réalisée aux besoins formulés au chapitre~\ref{chap:conception}. Il présente le plan de tests, les résultats de sécurité et de performance, les métriques du module ML, puis discute les limites et les perspectives.

\section{Plan de tests}
\label{sec:tc-plan}

Quatre familles de tests ont été définies, conformément aux besoins du chapitre~\ref{chap:conception}~:

\begin{table}[!htbp]
\centering
\caption{Plan de tests}
\label{tab:plan-tests}
\small
\begin{tabularx}{\linewidth}{|l|X|X|}
\hline
\textbf{Famille} & \textbf{Objectif} & \textbf{Outils} \\
\hline
Tests unitaires & Vérifier scripts Vault, simulateur, ETL ML & pytest, scripts shell \\ \hline
Tests d'intégration & Valider la chaîne PKI~$\to$~mTLS~$\to$~broker~$\to$~Suricata~$\to$~Wazuh & mosquitto\_pub/sub, eve.json \\ \hline
Tests de sécurité & Rejouer les six scénarios d'attaque & flotte simulée, scripts attaque \\ \hline
Tests de performance & Mesurer latence, débit, handshake mTLS & paho-mqtt, hyperfine \\ \hline
\end{tabularx}
\end{table}

\subsection{Tests unitaires et vérifications préalables}
\label{subsec:tc-unitaires}

Avant les campagnes de test, un ensemble de vérifications unitaires a validé chaque brique~: émission, renouvellement et révocation PKI Vault avec vérification de la chaîne de confiance (\texttt{openssl verify})~; connexion mTLS réussie avec certificat valide et rejet avec certificat expiré, auto-signé ou révoqué~; enforcement ACL Mosquitto (publication autorisée sur les topics du CN, rejet sur ceux d'un autre dispositif)~; chargement des trois modèles \texttt{.pkl} et prédiction sur un vecteur de test connu~; validation du routage inter-VLAN et du port mirroring DMZ\textrightarrow{}Suricata. Ces vérifications ont permis de corriger plusieurs erreurs de configuration avant le lancement des campagnes.

\section{Résultats de sécurité}
\label{sec:tc-secu}

Les six scénarios d'attaque définis en \S\ref{sec:re-attaques} ont été rejoués 10~fois chacun. Les taux de détection sont consignés dans le tableau~\ref{tab:detection-secu}.

\begin{table}[!htbp]
\centering
\caption{Taux de détection des scénarios d'attaque (10 rejeux par scénario)}
\label{tab:detection-secu}
\begin{tabularx}{\linewidth}{|X|c|c|c|}
\hline
\textbf{Scénario} & \textbf{Suricata} & \textbf{Wazuh (corr.)} & \textbf{ML} \\
\hline
S1 -- MITM TLS désactivé        & 10/10 & 10/10 & N/A \\ \hline
S2 -- Sniffing MQTT clair       & 10/10 & 10/10 & N/A \\ \hline
S3 -- Brute-force CONNECT       & 10/10 & 10/10 & 9/10 \\ \hline
S4 -- Topic non autorisé        &  8/10 & 10/10 & 10/10 \\ \hline
S5 -- Flood CONNECT             & 10/10 & 10/10 & 9/10 \\ \hline
S6 -- Compromission certificat  &  3/10 &  6/10 & 10/10 \\ \hline
\end{tabularx}
\end{table}

\begin{resultbox}[Lecture]
La complémentarité \emph{règles}~$+$~\emph{ML} est clairement démontrée~: les attaques exploitant un vecteur réseau identifiable (S1--S5) sont couvertes par les règles~; le scénario~S6 (certificat valide réutilisé) n'est détecté de manière fiable que par le module ML, qui repère l'écart comportemental.
\end{resultbox}

\subsection{Preuves visuelles}

Les captures ci-dessous étayent les taux du tableau~\ref{tab:detection-secu}.

\begin{figure}[!htbp]
\centering
\includegraphics[width=\textwidth]{atk-pfsense-rules.png}
\caption{pfSense~-- règles de filtrage inter-VLAN actives validant la politique \emph{deny-by-default} (O1)}
\label{fig:pfsense-rules}
\end{figure}

\begin{figure}[!htbp]
\centering
\includegraphics[width=\textwidth]{atk-fleet-sim.png}
\caption{Campagne complémentaire~: 100~dispositifs actifs, 35\,085~messages MQTT (O4)}
\label{fig:fleet-sim}
\end{figure}

\begin{figure}[!htbp]
\centering
\includegraphics[width=\textwidth]{atk-mosquitto-acl.png}
\caption{Logs Mosquitto~: connexions mTLS identifiées par CN et refus ACL sur topics non autorisés (O3, BF-3)}
\label{fig:mosquitto-acl}
\end{figure}
\FloatBarrier

La figure~\ref{fig:pfsense-rules} confirme la politique \emph{deny-by-default}~: seuls IoT\textrightarrow{}Broker/8883 et IoT\textrightarrow{}Vault/8200 sont autorisés. La figure~\ref{fig:fleet-sim} valide la montée en charge à 100~dispositifs (O4). La figure~\ref{fig:mosquitto-acl} montre les connexions mTLS par CN et les refus ACL.

\begin{figure}[!htbp]
\centering
\includegraphics[width=\textwidth]{atk-flood-summary.png}
\caption{Scénario S5 -- flood CONNECT~: 24~tentatives, 0~connexion réussie, \texttt{exit code 1}}
\label{fig:atk-summary}
\end{figure}

\begin{figure}[!htbp]
\centering
\includegraphics[width=\textwidth]{suricata-alerts.png}
\caption{Alertes Suricata~: SID 9000001--9000004, 51~alertes, 0~faux positif sur trafic TLS (O5)}
\label{fig:suricata-alerts-test}
\end{figure}

\begin{figure}[!htbp]
\centering
\includegraphics[width=\textwidth]{wazuh-dashboard.png}
\caption{Tableau de bord Wazuh~: alertes corrélées par criticité et type d'attaque (O6, BNF-1)}
\label{fig:wazuh-dashboard-test}
\end{figure}
\FloatBarrier

Le blocage du flood S5 (figure~\ref{fig:atk-summary}) confirme l'efficacité du filtrage pfSense + Mosquitto. Les figures~\ref{fig:suricata-alerts-test} et~\ref{fig:wazuh-dashboard-test} valident la chaîne de détection Suricata\textrightarrow{}Wazuh de bout en bout.

\section{Résultats du module ML}
\label{sec:tc-ml}

\subsection{Isolation Forest}

\begin{figure}[!htbp]
\centering
\includegraphics[width=0.88\linewidth]{iso_forest_confusion_matrix.png}
\caption{Isolation Forest -- matrice de confusion binaire}
\label{fig:iso-eval}
\end{figure}

\begin{figure}[!htbp]
\centering
\includegraphics[width=0.88\linewidth]{iso_forest_scores_distribution.png}
\caption{Isolation Forest -- distribution des scores d'anomalie (seuil à $-0{,}05$)}
\label{fig:iso-scores}
\end{figure}
\FloatBarrier

Le modèle distingue correctement le trafic normal des anomalies avec un rappel de 0{,}947. La distribution des scores confirme la séparabilité des deux populations.

\subsection{Random Forest}

\begin{figure}[!htbp]
\centering
\includegraphics[width=0.95\linewidth]{rf_confusion_roc.png}
\caption{Random Forest~: matrice de confusion binaire et courbe ROC (AUC~=~0{,}999)}
\label{fig:rf-eval}
\end{figure}

\begin{figure}[!htbp]
\centering
\includegraphics[width=0.88\linewidth]{rf_multiclass_confusion.png}
\caption{Random Forest -- confusion multi-classes~: 7~types d'attaques, absence de confusion Normal$\leftrightarrow$DoS}
\label{fig:rf-multiclass}
\end{figure}

\begin{figure}[!htbp]
\centering
\includegraphics[width=0.90\linewidth]{rf_feature_importance.png}
\caption{Importance des features~: \emph{inter-arrival time} et \emph{payload length} sont les plus discriminants}
\label{fig:rf-features}
\end{figure}
\FloatBarrier

\begin{table}[!htbp]
\centering
\caption{Performances ML (validation 5-fold)}
\label{tab:perf-ml}
\begin{tabularx}{\linewidth}{|X|c|c|c|c|}
\hline
\textbf{Modèle} & \textbf{Accuracy} & \textbf{Precision} & \textbf{Recall} & \textbf{F1} \\
\hline
Isolation Forest (binaire)  & 0{,}962 & 0{,}918 & 0{,}947 & 0{,}932 \\ \hline
Random Forest (multi-classe)& 0{,}994 & 0{,}993 & 0{,}993 & 0{,}993 \\ \hline
\end{tabularx}
\end{table}

\begin{resultbox}[Interprétation analytique]
L'écart de F1 (0{,}932 vs 0{,}993) reflète la différence de paradigmes~: IF modélise la normalité sans étiquettes et couvre les anomalies inconnues~; RF bénéficie des étiquettes et maximise la précision. Les features les plus discriminantes (\emph{inter-arrival time}, longueur payload, figure~\ref{fig:rf-features}) confirment les signatures volumétriques des attaques. L'usage conjoint est validé~: IF détecte S6, RF classifie S3--S5.
\end{resultbox}

\section{Résultats de performance}
\label{sec:tc-perf}

\begin{figure}[!htbp]
\centering
\includegraphics[width=0.88\linewidth]{perf_tls_handshake.png}
\caption{Distribution du handshake mTLS}
\end{figure}

\begin{figure}[!htbp]
\centering
\includegraphics[width=0.88\linewidth]{perf_connection_latency.png}
\caption{Latence de connexion}
\end{figure}
\FloatBarrier

Le handshake mTLS médian s'établit à $\sim 38$~ms (95\textsuperscript{e} percentile~: 47~ms), conforme à BNF-2 (<~50~ms).

\begin{figure}[!htbp]
\centering
\includegraphics[width=0.88\linewidth]{perf_throughput.png}
\caption{Débit applicatif}
\end{figure}

\begin{figure}[!htbp]
\centering
\includegraphics[width=0.88\linewidth]{perf_rtt_messages.png}
\caption{RTT par taille de message}
\end{figure}

\begin{figure}[!htbp]
\centering
\includegraphics[width=0.85\linewidth]{perf_protocols_comparison.png}
\caption{Comparaison MQTT clair vs MQTT/TLS vs MQTT/mTLS}
\end{figure}
\FloatBarrier

\begin{resultbox}[Validation BNF-2]
Le surcoût mTLS représente $\sim 8$\,\% du débit et $\sim 3{,}5$~ms de latence supplémentaire. Le débit agrégé ($\sim 31\,000$~msg/min, soit $\sim 620$~msg/s en pic et $\sim 520$~msg/s soutenu) dépasse l'exigence de 500~msg/s par client. Le surcoût ne compromet pas la capacité du broker.
\end{resultbox}

\section{Limites identifiées}
\label{sec:tc-limites}

\begin{itemize}[leftmargin=*]
  \item \textbf{Environnement simulé}~: les mesures ne reflètent pas les conditions radio réelles (latence, perte de paquets, mobilité).
  \item \textbf{Scalabilité partielle}~: le périmètre nominal (50~dispositifs) et la campagne complémentaire (100~dispositifs) valident les premiers paliers. Le palier BNF-4 (200~dispositifs) reste à vérifier~; la contrainte est matérielle (16~Go RAM), non architecturale.
  \item \textbf{ML hors ligne}~: l'intégration en streaming (Kafka + scoring temps réel) reste à explorer.
  \item \textbf{Rotation de certificats}~: une intégration complète au cycle ACME serait souhaitable.
\end{itemize}

\section{Perspectives}
\label{sec:tc-perspectives}

\textbf{Court terme} ($<$~3~mois)~: industrialisation des scripts (Ansible), intégration ACME/Vault, extension à 200~clients.\\
\textbf{Moyen terme} (3--12~mois)~: portage cellulaire réel (4G/5G), intégration aux outils SOC de \TT{}, entraînement ML continu.\\
\textbf{Long terme} ($>$~12~mois)~: architecture \emph{Zero Trust} IoT, intégration LoRaWAN/NB-IoT, déploiement multi-tenant cloud.

\section{Conclusion}
\label{sec:tc-synthese}

Les résultats confirment la validité de l'architecture proposée. Les sept objectifs (O1--O7) sont atteints~: segmentation opérationnelle, PKI automatisée, mTLS obligatoire, flotte simulée (campagne complémentaire à 100~dispositifs), chaîne Suricata--Wazuh et module ML avec F1~=~0{,}993 en multi-classes. Le surcoût mTLS reste modéré ($\sim 8$\,\%). Les limites identifiées définissent le périmètre de validité et orientent les perspectives vers une solution de grade opérationnel.
